\documentclass[11pt]{article}

\usepackage[margin=1in]{geometry}
\usepackage{amsmath,amssymb,amsthm}
\usepackage{enumitem}
\usepackage{xcolor}
\usepackage{hyperref}

\hypersetup{
  colorlinks=true,
  linkcolor=blue!60!black,
  urlcolor=blue!60!black
}

\newcommand{\R}{\mathbb{R}}
\newcommand{\N}{\mathbb{N}}
\newcommand{\F}{\mathcal{F}}
\newcommand{\Prob}{\mathbb{P}}
\newcommand{\E}{\mathbb{E}}
\newcommand{\dd}{\,d}
\newcommand{\Knowledge}[1]{%
  \par\smallskip\noindent
  \textbf{Knowledge used.}
  {\small\color{blue!60!black}\raggedright #1\par}
  \par\smallskip
}

\title{Chapter 8 Exercises}
\author{\textsc{Solutions to Rick Durrett's}\\
\textit{Probability: Theory and Examples}}
\date{\today}

\begin{document}
\maketitle

\noindent
These solutions use the local Chapter 8 LLM/Viki knowledge set in
\path{Probability/LLM_Viki_Dataset/Chapter8_Knowledge}.  The cited
knowledge IDs refer to entries in
\path{chapter8_knowledge_pieces.jsonl}.

\section*{Exercise 8.1.1}

\noindent\textbf{Question.}
Use Exercise 7.5.4 to conclude that
\[
  \E T_{U,V}^2 \le 4\E X^4.
\]

\Knowledge{
\path{durrett_ch8_skorokhod_embedding_mean_zero};
\path{durrett_ch8_embedded_random_walk}.
}

\noindent\textbf{Solution.}
Recall the notation from the proof of Skorokhod's representation
theorem.  The mean-zero law of $X$ is written as a mixture of two-point
mean-zero laws $\mu_{u,v}$, where $u<0<v$ and
\[
  \mu_{u,v}(\{u\})=\frac{v}{v-u},
  \qquad
  \mu_{u,v}(\{v\})=\frac{-u}{v-u}.
\]
Given $(U,V)=(u,v)$, the stopping time is
\[
  T_{u,v}=\inf\{t:B_t\notin(u,v)\},
\]
so $B_{T_{u,v}}$ has law $\mu_{u,v}$.

Exercise 7.5.4 says that for Brownian motion exiting an interval
$(a,b)$ with $a<0<b$,
\[
  \E T_{a,b}^2 \le 4 \E B_{T_{a,b}}^4.
\]
Applying this conditional on $(U,V)=(u,v)$ gives
\[
  \E\!\left(T_{U,V}^2\mid U,V\right)
  \le
  4\E\!\left(B_{T_{U,V}}^4\mid U,V\right)
  =
  4\int x^4\,\mu_{U,V}(\dd x).
\]
Now take expectations over the mixing distribution of $(U,V)$.  Since
the mixture of the laws $\mu_{U,V}$ is exactly the original law of $X$,
\[
  \E T_{U,V}^2
  \le
  4 \E\int x^4\,\mu_{U,V}(\dd x)
  =
  4\E X^4.
\]
If $\E X^4=\infty$, the inequality is still true in the extended sense;
otherwise it gives the desired finite second-moment bound for the
embedding time.

\section*{Exercise 8.1.2}

\noindent\textbf{Question.}
Suppose $S_n$ is one-dimensional simple random walk and let
\[
  R_n
  =
  1+\max_{m\le n}S_m-\min_{m\le n}S_m
\]
be the number of points visited by time $n$.  Show that $R_n/\sqrt n$
converges in distribution to a limit.

\Knowledge{
\path{durrett_ch8_donsker_invariance_principle};
\path{durrett_ch8_continuous_functional_transfer};
\path{durrett_ch8_maximum_functional}.
}

\noindent\textbf{Solution.}
For one-dimensional nearest-neighbor simple random walk, every integer
between the running minimum and running maximum has been visited.  Hence
the number of visited points by time $n$ is exactly
\[
  R_n
  =
  1+\max_{0\le m\le n}S_m-\min_{0\le m\le n}S_m.
\]

Let $W_n(t)=S(\lfloor nt\rfloor)/\sqrt n$ with linear interpolation
between the mesh points.  Since the increments of simple random walk
have mean $0$ and variance $1$, Donsker's theorem gives
\[
  W_n \Rightarrow B
  \qquad\text{in } C[0,1],
\]
where $B$ is standard Brownian motion.

Define the range functional
\[
  \psi(\omega)=\sup_{0\le t\le 1}\omega(t)
  -\inf_{0\le t\le 1}\omega(t),
  \qquad \omega\in C[0,1].
\]
This functional is continuous under the uniform norm.  Indeed, if
$\|\omega-\eta\|_\infty\le\varepsilon$, then
\[
  |\sup_t\omega(t)-\sup_t\eta(t)|\le\varepsilon,
  \qquad
  |\inf_t\omega(t)-\inf_t\eta(t)|\le\varepsilon,
\]
and therefore $|\psi(\omega)-\psi(\eta)|\le 2\varepsilon$.
The continuous mapping theorem now implies
\[
  \max_{0\le m\le n}\frac{S_m}{\sqrt n}
  -
  \min_{0\le m\le n}\frac{S_m}{\sqrt n}
  =
  \psi(W_n)
  \Rightarrow
  \psi(B)
  =
  \sup_{0\le t\le 1}B_t-\inf_{0\le t\le 1}B_t.
\]
Finally, $1/\sqrt n\to0$, so Slutsky's theorem gives
\[
  \frac{R_n}{\sqrt n}
  =
  \frac{1}{\sqrt n}
  +
  \max_{0\le m\le n}\frac{S_m}{\sqrt n}
  -
  \min_{0\le m\le n}\frac{S_m}{\sqrt n}
  \Rightarrow
  \sup_{0\le t\le 1}B_t-\inf_{0\le t\le 1}B_t.
\]
Thus $R_n/\sqrt n$ converges in distribution.  The limiting random
variable is the range of Brownian motion on $[0,1]$.

\section*{Section 8.2 Exercises}

\noindent\textbf{Question.}
The local textbook PDF has no separately numbered exercises in
Section 8.2.  Section 8.2, ``CLTs for Martingales,'' begins immediately
after Exercise 8.1.2 and then proceeds through Theorems 8.2.1--8.2.8
before Section 8.3 begins.

\Knowledge{
\path{durrett_ch8_martingale_sk_embedding};
\path{durrett_ch8_martingale_convergence_variance_sum};
\path{durrett_ch8_martingale_fclt_random_clock};
\path{durrett_ch8_martingale_lindeberg_feller};
\path{durrett_ch8_martingale_clt_final_form}.
}

\noindent\textbf{Solution.}
There are therefore no Section 8.2 exercise statements to solve from
this source file.  The main usable knowledge from the section is the
martingale CLT toolkit:
\begin{enumerate}[label=(\roman*)]
\item square-integrable martingales can be embedded in Brownian motion
  at stopping times;
\item the Brownian clock is controlled by conditional variances;
\item a martingale functional CLT follows when the conditional variance
  process converges to deterministic time and large jumps are negligible;
\item the Lindeberg-Feller martingale theorem is obtained by truncation;
\item the single-sequence martingale CLT follows by converting the
  sequence into a triangular array.
\end{enumerate}
These are recorded above as the relevant Viki knowledge IDs, but the
textbook does not give an Exercise 8.2.1 or any other Section 8.2
exercise to solve.

\section*{Exercise 8.3.1}

\noindent\textbf{Question.}
Consider the special case of Example 8.3.3 in which $\xi_i$, $i\in\mathbb Z$,
are i.i.d. and take values $H$ and $T$ with equal probability.  Let
\[
  \eta_n=f(\xi_n,\xi_{n+1}),
  \qquad
  f(H,T)=1,
  \qquad
  f(i,j)=0 \text{ otherwise}.
\]
Then $S_n=\eta_1+\cdots+\eta_n$ counts the number of head runs in
$\xi_1,\ldots,\xi_{n+1}$.  Apply Theorem 8.3.2 with
$\F_n=\sigma(\xi_m:m\le n+1)$ to show that there are constants $\mu$
and $\sigma$ so that
\[
  \frac{S_n-n\mu}{\sigma n^{1/2}}\Rightarrow \chi .
\]
What is the random variable $Y_0$ constructed in the proof of
Theorem 8.3.2 in this case?

\Knowledge{
\path{durrett_ch8_stationary_coboundary_clt};
\path{durrett_ch8_projective_stationary_clt};
\path{durrett_ch8_m_dependent_clt};
\path{durrett_ch8_martingale_clt_final_form}.
}

\noindent\textbf{Solution.}
Let
\[
  X_n=\eta_n-\mu,
  \qquad
  \mu=\E\eta_n.
\]
Since $\eta_n=1_{\{\xi_n=H,\xi_{n+1}=T\}}$,
\[
  \mu=\Prob(\xi_n=H,\xi_{n+1}=T)=\frac14.
\]
The sequence $(X_n)$ is stationary and ergodic because it is a
shift-factor of the i.i.d. sequence $(\xi_n)$.  Also $X_0$ depends only
on $\xi_0,\xi_1$.  With $\F_n=\sigma(\xi_m:m\le n+1)$, for $k\ge2$ the
$\sigma$-field $\F_{-k}$ is generated by variables $\xi_m$ with
$m\le -k+1\le -1$, and hence is independent of $X_0$.  Therefore
\[
  \E(X_0\mid \F_{-k})=0,
  \qquad k\ge2.
\]
Thus
\[
  \sum_{k\ge1}\|\E(X_0\mid \F_{-k})\|_2<\infty,
\]
so Theorem 8.3.2 applies.

It remains to compute the constants.  We already have $\mu=1/4$.  Since
$(\eta_n)$ is $1$-dependent, only the lag-one covariance contributes:
\[
  \sigma^2
  =
  \operatorname{Var}(\eta_0)
  +2\operatorname{Cov}(\eta_0,\eta_1).
\]
Now
\[
  \operatorname{Var}(\eta_0)
  =
  \frac14\left(1-\frac14\right)
  =
  \frac{3}{16}.
\]
The events $\{\eta_0=1\}$ and $\{\eta_1=1\}$ cannot occur together,
because they would require $\xi_1=T$ and $\xi_1=H$ simultaneously.
Hence $\E(\eta_0\eta_1)=0$, and
\[
  \operatorname{Cov}(\eta_0,\eta_1)
  =
  0-\frac14\cdot\frac14
  =
  -\frac1{16}.
\]
Therefore
\[
  \sigma^2=\frac{3}{16}+2\left(-\frac1{16}\right)=\frac1{16},
  \qquad
  \sigma=\frac14.
\]
Theorem 8.3.2 gives
\[
  \frac{S_n-n/4}{(1/4)\sqrt n}\Rightarrow \chi,
\]
where $\chi$ has the standard normal distribution.

Finally we identify the $Y_0$ from the proof of Theorem 8.3.2.  Write
$I_j=1_{\{\xi_j=H\}}$, so $\eta_j=I_j(1-I_{j+1})$ and
$X_j=\eta_j-1/4$.  The proof defines
\[
  Y_0=\sum_{j=0}^\infty
  \left\{\E(X_j\mid \F_0)-\E(X_j\mid \F_{-1})\right\}.
\]
All terms with $j\ge2$ vanish, because $X_j$ is independent of $\F_0$.
For $j=0$,
\[
  \E(X_0\mid \F_0)=X_0,
  \qquad
  \E(X_0\mid \F_{-1})=\frac12 I_0-\frac14.
\]
For $j=1$,
\[
  \E(X_1\mid \F_0)=\frac12 I_1-\frac14,
  \qquad
  \E(X_1\mid \F_{-1})=0.
\]
Thus
\[
  Y_0
  =
  \left(I_0(1-I_1)-\frac14\right)
  -\left(\frac12 I_0-\frac14\right)
  +\left(\frac12 I_1-\frac14\right),
\]
or equivalently
\[
  Y_0
  =
  \frac12 I_0+\frac12 I_1-I_0I_1-\frac14.
\]
This takes the value $1/4$ when $\xi_0\ne \xi_1$ and $-1/4$ when
$\xi_0=\xi_1$.  In particular $\E Y_0^2=1/16$, matching the limiting
variance computed above.

\section*{Exercise 8.4.1}

\noindent\textbf{Question.}
Use Exercise 7.4.8 and the reasoning that led to (8.4.11) to conclude
\[
  \Prob\left(\max_{0\le t\le1}B_t^0>b\right)=\exp(-2b^2),
\]
where $B^0$ is Brownian bridge.

\Knowledge{
\path{durrett_ch8_brownian_bridge_conditioning};
\path{durrett_ch8_kolmogorov_distribution_method}.
}

\noindent\textbf{Solution.}
Think of Brownian bridge as Brownian motion conditioned on $B_1=0$.
Let $T_b=\inf\{t:B_t=b\}$.  For Brownian motion starting at $0$,
\[
  \left\{\max_{0\le t\le1}B_t>b\right\}
  =
  \{T_b<1\}.
\]
Exercise 7.4.8 gives the reflected density
\[
  \Prob_0(T_b<1,\ B_1\in dx)
  =
  p_1(0,2b-x)\,dx,
  \qquad x<b.
\]
Conditioning on $B_1=0$ therefore gives
\[
  \Prob\left(\max_{0\le t\le1}B_t^0>b\right)
  =
  \frac{p_1(0,2b)}{p_1(0,0)}.
\]
Since $p_1(0,x)=(2\pi)^{-1/2}\exp(-x^2/2)$,
\[
  \frac{p_1(0,2b)}{p_1(0,0)}
  =
  \exp\left(-\frac{(2b)^2}{2}\right)
  =
  \exp(-2b^2).
\]

\section*{Exercise 8.4.2}

\noindent\textbf{Question.}
Let $B_t$ be a Brownian motion starting at $0$.  Show that
\[
  X_t=(1-t)B\left(\frac{t}{1-t}\right), \qquad 0\le t<1,
\]
is a Brownian bridge.

\Knowledge{
\path{durrett_ch8_brownian_bridge_conditioning}.
}

\noindent\textbf{Solution.}
The process $X$ is Gaussian and has mean zero, since it is obtained by
deterministic linear transformations of Brownian motion.  For
$0\le s\le t<1$,
\[
  \E(X_sX_t)
  =
  (1-s)(1-t)
  \E\left[
    B\left(\frac{s}{1-s}\right)
    B\left(\frac{t}{1-t}\right)
  \right].
\]
Because $u\mapsto u/(1-u)$ is increasing,
\[
  \E\left[
    B\left(\frac{s}{1-s}\right)
    B\left(\frac{t}{1-t}\right)
  \right]
  =
  \frac{s}{1-s}.
\]
Thus
\[
  \E(X_sX_t)
  =
  (1-s)(1-t)\frac{s}{1-s}
  =
  s(1-t),
\]
which is exactly the covariance of Brownian bridge
$B_t^0=B_t-tB_1$.

Finally, as $t\uparrow1$ let $u=t/(1-t)$.  Then
\[
  X_t=\frac{B_u}{1+u}.
\]
Brownian motion is sublinear at infinity, so $B_u/(1+u)\to0$ almost
surely.  Setting $X_1=0$, the process has continuous paths on $[0,1]$.
Since it is a centered Gaussian process with Brownian-bridge covariance,
$X$ is a Brownian bridge.

\section*{Exercise 8.4.3}

\noindent\textbf{Question.}
Show that Brownian bridge is a Markov process by checking that if
$s_1<\cdots<s_n<s<t<1$, then
\[
  \Prob(B^0(t)=y\mid B^0(s)=x,B^0(s_n)=x_n,\ldots,B^0(s_1)=x_1)
  =
  \Prob(B^0(t)=y\mid B^0(s)=x).
\]

\Knowledge{
\path{durrett_ch8_brownian_bridge_conditioning};
\path{durrett_ch8_kolmogorov_distribution_method}.
}

\noindent\textbf{Solution.}
Use the pinned Brownian transition-density formula.  If
$0=t_0<s_1<\cdots<s_n<s<t<1$ and $x_0=0$, the joint density of
\[
  B^0(s_1),\ldots,B^0(s_n),B^0(s),B^0(t)
\]
at $x_1,\ldots,x_n,x,y$ is proportional to
\[
  \left\{\prod_{j=1}^n
  p_{s_j-s_{j-1}}(x_{j-1},x_j)\right\}
  p_{s-s_n}(x_n,x)\,
  p_{t-s}(x,y)\,
  p_{1-t}(y,0).
\]
The joint density of
$B^0(s_1),\ldots,B^0(s_n),B^0(s)$ at
$x_1,\ldots,x_n,x$ is proportional to
\[
  \left\{\prod_{j=1}^n
  p_{s_j-s_{j-1}}(x_{j-1},x_j)\right\}
  p_{s-s_n}(x_n,x)\,
  p_{1-s}(x,0).
\]
Dividing the two densities, all factors involving the earlier times
$s_1,\ldots,s_n$ cancel.  Hence the conditional density of $B^0(t)$
given the displayed past is
\[
  \frac{p_{t-s}(x,y)p_{1-t}(y,0)}{p_{1-s}(x,0)},
\]
which depends only on the present value $B^0(s)=x$.  This proves the
Markov property.

\section*{Exercise 8.4.4}

\noindent\textbf{Question.}
Let $B_t^0$ be Brownian bridge.  Show that
\[
  X_t=\frac{B_t^0}{1-t}, \qquad 0\le t<1,
\]
is a martingale but is not $L^2$ bounded.

\Knowledge{
\path{durrett_ch8_brownian_bridge_conditioning}.
}

\noindent\textbf{Solution.}
Let $\F_t^0=\sigma(B_u^0:0\le u\le t)$ be the natural filtration of the
bridge.  From the Brownian-bridge transition law, for $0\le s<t<1$,
\[
  \E(B_t^0\mid \F_s^0)
  =
  \E(B_t^0\mid B_s^0)
  =
  \frac{1-t}{1-s}B_s^0.
\]
Therefore
\[
  \E(X_t\mid \F_s^0)
  =
  \frac{1}{1-t}\E(B_t^0\mid \F_s^0)
  =
  \frac{B_s^0}{1-s}
  =
  X_s.
\]
Thus $(X_t)_{0\le t<1}$ is a martingale.

It is not $L^2$ bounded.  Since Brownian bridge has
\[
  \operatorname{Var}(B_t^0)=t(1-t),
\]
we have
\[
  \E X_t^2
  =
  \frac{\E(B_t^0)^2}{(1-t)^2}
  =
  \frac{t(1-t)}{(1-t)^2}
  =
  \frac{t}{1-t}.
\]
As $t\uparrow1$, this tends to infinity.  Hence
\[
  \sup_{0\le t<1}\E X_t^2=\infty,
\]
so the martingale is not bounded in $L^2$.

\section*{Exercise 8.5.1}

\noindent\textbf{Question.}
Show that if $\E|X_i|^\alpha=\infty$ for some $\alpha<2$, then
\[
  \limsup_{n\to\infty}\frac{|X_n|}{n^{1/\alpha}}=\infty
  \qquad\text{a.s.}
\]
so the law of the iterated logarithm fails.

\Knowledge{
\path{durrett_ch8_iid_lil};
\path{durrett_ch8_brownian_lil}.
}

\noindent\textbf{Solution.}
Let $Y=|X_1|^\alpha$.  Since $\E Y=\infty$ and the tail function
$t\mapsto\Prob(Y>t)$ is nonincreasing,
\[
  \sum_{n=1}^\infty \Prob(Y>cn)=\infty
  \qquad\text{for every } c>0.
\]
Equivalently, for every $M>0$,
\[
  \sum_{n=1}^\infty
  \Prob\left(|X_n|>M n^{1/\alpha}\right)
  =
  \sum_{n=1}^\infty
  \Prob\left(Y>M^\alpha n\right)
  =
  \infty.
\]
The events $\{|X_n|>M n^{1/\alpha}\}$ are independent, so the second
Borel-Cantelli lemma gives
\[
  |X_n|>M n^{1/\alpha}
  \quad\text{infinitely often a.s.}
\]
Since $M$ is arbitrary, applying this to $M=1,2,3,\ldots$ gives
\[
  \limsup_{n\to\infty}\frac{|X_n|}{n^{1/\alpha}}=\infty
  \qquad\text{a.s.}
\]

This rules out the LIL normalization.  Indeed, if an increment is very
large, then at least one of the two adjacent partial sums is large:
\[
  |X_n|=|S_n-S_{n-1}|
  \le |S_n|+|S_{n-1}|,
\]
so
\[
  \max(|S_n|,|S_{n-1}|)\ge \frac{|X_n|}{2}.
\]
Since $\alpha<2$,
\[
  \frac{n^{1/\alpha}}{\sqrt{n\log\log n}}
  =
  \frac{n^{1/\alpha-1/2}}{\sqrt{\log\log n}}
  \longrightarrow \infty.
\]
Along the infinitely many large-jump indices above, the LIL scale
$\sqrt{2n\log\log n}$ is therefore too small to keep the partial sums
bounded.  Hence the finite-variance LIL conclusion cannot hold.

\section*{Exercise 8.5.2}

\noindent\textbf{Question.}
Give a direct proof that, under the hypotheses of Theorem 8.5.3, the
limit set of
\[
  \left\{
  \frac{S_n}{(2n\log\log n)^{1/2}}: n\ge3
  \right\}
\]
is $[-1,1]$.

\Knowledge{
\path{durrett_ch8_iid_lil};
\path{durrett_ch8_strassen_invariance_principle}.
}

\noindent\textbf{Solution.}
Theorem 8.5.3 says that the limit set in $C[0,1]$ of
\[
  Z_n(t)=\frac{S(nt)}{(2n\log\log n)^{1/2}}
\]
is
\[
  K=\left\{
  f:f(x)=\int_0^x g(y)\,dy,\ 
  \int_0^1 g(y)^2\,dy\le1
  \right\}.
\]
The map $e_1:C[0,1]\to\R$ given by $e_1(f)=f(1)$ is continuous.
Therefore the scalar subsequential limits of
\[
  \frac{S_n}{(2n\log\log n)^{1/2}}=Z_n(1)
\]
are exactly the values $f(1)$ with $f\in K$.

First let $f\in K$.  Then $f(1)=\int_0^1 g(y)\,dy$ for some $g$ with
$\int_0^1g^2\le1$.  By Cauchy's inequality,
\[
  |f(1)|
  =
  \left|\int_0^1 g(y)\,dy\right|
  \le
  \left(\int_0^1g(y)^2\,dy\right)^{1/2}
  \le1.
\]
Thus every scalar limit point lies in $[-1,1]$.

Conversely, fix $a\in[-1,1]$ and set $f_a(t)=at$.  Then
$f_a(t)=\int_0^t a\,dy$ and
\[
  \int_0^1 a^2\,dy=a^2\le1,
\]
so $f_a\in K$.  By Theorem 8.5.3, there is a subsequence $n_k$ such that
$Z_{n_k}\to f_a$ uniformly on $[0,1]$.  Evaluating at $1$ gives
\[
  \frac{S_{n_k}}{(2n_k\log\log n_k)^{1/2}}
  =
  Z_{n_k}(1)
  \longrightarrow
  f_a(1)=a.
\]
Hence every $a\in[-1,1]$ is a limit point, and the scalar limit set is
exactly $[-1,1]$.

\end{document}
